\section{The City of the Sun}

The \textbf{City of the Sun: A Poetic Dialog} is a Utopian plan devised by the Hermetist and Dominican monk, \textbf{Thomas Campanella} to describe the ideal Solar civilization. Due to being on the wrong side of political intrigue and his alleged heresy, Campanella spent 27 years in prison, where he wrote most of his works. Despite this, he amazingly held an optimistic attitude, even for the Spanish monarchy and Roman church which were his persecutors.

The novel is presented as the report of an adventurer to questions asked by a Knight on his return. The adventurer was forced to land on an Island in the South Seas, where he stumbled upon the isolated City of the Sun, known by the alchemical symbol for the sun.

\paragraph{Rule}


The City is ruled by the Metaphysician (or Sun), who serves as both spiritual and temporal leader. Beneath him are three assistants:

\begin{itemize}


\item \textbf{Power}, in charge of military affairs


\item \textbf{Wisdom}, in charge of science, liberal arts and technology


\item \textbf{Love}, in charge of breeding, agriculture, medicine, and education
\end{itemize}


Beneath these three, are specialists in each of the areas they rule over. So the Metaphysician is not only the expert in spiritual matters, but he must also have mastered the empirical sciences as well as administrative and practical techniques.

The obvious question does come up, since philosophers and contemplatives don't seem to also have the required leadership qualities. The answer is that the Sun is much more knowledgeable than what passes for knowledge in Europe and furthermore the European system of choosing rulers on the basis of birth or power factions is hardly a wise system. We would add that, in our day, the system of democracy, or appeal to mass popularity, is not a wise system for selecting rulers.

Evola solves the problem from a different angle. For him, the Emperor also takes a sacerdotal initiation, which makes him the head of the priestly caste as well as the warrior caste. Thus he already has temporal power and then adds the spiritual authority to it.

Both schemes differ from the thought of Rene Guenon and the Roman Church, who hold to a dual leadership plan: the spiritual authority (Pope or \emph{Brahmin}) and temporal power (Emperor or
\emph{Kshatriya}). Evola is highly critical of the dual scheme, which he calls ``two suns'', an impossibility in his mind.

Thus both Evola and Campanella advocate forms of theocracy. The Guenonian position may be described as theonomy, with the spiritual authority setting the principles and the temporal power implementing them.

\paragraph{Religion}


As the inhabitants are originally from India, their religion derives from a purified Brahminism (or we would say the Vedanta), with the Sun representing the one God. Since they are knowledgeable of all the nations in the world, they admire the great thinkers and religious leaders (including Mohammed). Their highest regard is for Jesus Christ and the apostles, although, we must point out, they never converted to Christianity. This interest in the best minds in all of history anticipates Comte's religion of humanity by two centuries.

The Solarians believe in two principles: the sun as father and the earth as mother. These are obviously the \textbf{Solar} and \textbf{Lunar} principles.

Campanella describes it as a natural religion, in perfect harmony with Christianity, lacking only the sacraments. However, we must point out, for the Church the sacraments are essential to living a virtuous life, but apparently not for the City of the Sun. The implication is that the natural religion of the island produces a society superior to the Christian societies of Europe.

\paragraph{Lifestyle}


In the City, all goods are held in common and everyone's needs are taken care of. Meals and lodging are also in common, and people are seldom alone. Lodgings are periodically rotated. Utopians are usually attracted to this sort of communism as the way to mitigate against the problems arising from material inequality. Since the City is so successful economically, there is no need for private property and everyone is taken care of. Everyone has his assigned work, so there is no unemployment. Nevertheless, there is ample time for leisure, which is spent in athletic games, since the workday is just four hours.

There is near equality between the sexes, with the exception that women are not in leadership positions and are exempt from certain physical tasks. They even participate in the military.

\paragraph{Breeding}


In the City, women are also held in common, and permanent male-female relationships don't exist. Mating is determined by the authorities who match up couples based on physical characteristics and astrological signs. The latter implies that the citizens are all born around the same time in any given year, rather than randomly throughout the year. Those unfit to breed are offered sexual release either with barren (i.e., post-menopausal) or pregnant women. Children are raised in common after they are weaned (around two years old).

The City justifies this with the comment that Europeans are more careful with breeding horses and dogs than their own peoples. The problem with this logic is that we don't really breed the perfect horse or dog; rather, we create breeds of those animals geared for a specific purpose. Similarly, we can't say there is one ``ideal'' human type, since a properly functioning society needs people of various skills and interests. If we were to breed humans with as many varieties as we have for dogs, it is hard to imagine what sort of society would result; however, it would certainly be dysfunctional.

Nevertheless, in a limited sense, it has value. One is to maintain the racial purity of the City and the second to eliminate congenital diseases. It is hard to imagine that this assigned breeding would work well in practice, since there would be jealousy around who gets the most desirable mates. There are two other options we can mention.

The first is described by \textbf{Rousas Rushdoony} as the system in Armenia. There, he claims, the parents find suitable mates for their children based on genetic qualities. He claims that this has reduced congenital diseases in Armenia over the generations, though I am not aware of any empirical studies that back this up.

The other is the system of racial education proposed by Julius Evola, which he makes a part of general education. This counts on the citizenry themselves making wise mating choices. Note that all these schemes assume that children are produced for society, not just for the parents.

\paragraph{Education}


In the City, education is based on a complex system of pictorial representations which is apparently very effective in imparting education to the children.

\paragraph{Immigration and Trade}


Since the City is mostly self-sufficient, trade with foreigners is held to a minimum. However, the foreign merchants and the slaves acquired in warfare are not permitted to disrupt the orderliness of the City. A small number of immigrants are accepted, provided they demonstrate their willingness and ability to integrate into the lifestyle of the Solarians.

\paragraph{Conclusion}


Campanella concludes this work with the claim that the Solarians, based on their astrological calculations, have revealed to him the many major changes coming to the world. Campanella rightly points out that the 16th century produced more history and inventions than the preceding 4000 years; this trend is even accelerating in our time. One of his more interesting predictions is the increasing feminization of the West. He points to the increasing acceptance of pederasty, sodomy, and whoring among the literary elite. He claims, too, that ``men are becoming effeminate.'' He claims to know about other impending changes --- and without doubt there have been many since Campanella's time --- although he ends the book without revealing them. Nevertheless, whether because of the astrological influence of Venus or some other reason, the negative effects can be traced back to the \textbf{epicenity} of Western manhood.


\flrightit{Posted on 2010-06-06 by Cologero}

\begin{center}* * *\end{center}

\begin{footnotesize}
\begin{sffamily}
\texttt{Elsinore on 2010-06-10 at 00:40 said:}

Amazing article. It is no wonder that I so often frequent Gornahoor; my learning here is almost unparalleled to any other website on the net.

-Elsinore, Cordeliaforlear

\hfill

\texttt{Cologero on 2010-06-14 at 20:57 said:}

If we can get first principles established, the rest will follow.

\hfill

\texttt{kadambari on 2010-06-19 at 22:39 said:}

``As the inhabitants are originally from India, their religion derives from a purified Brahminism (or we would say the Vedanta), with the Sun representing the one God. Since they are knowledgeable of all the nations in the world, they admire the great thinkers and religious leaders (including Mohammed). Their highest regard is for Jesus Christ and the apostles, though, we must point out, they never converted to Christianity.''

As a ``Brahmin'' this made me laugh so much I nearly fell. Admire Mohammed? Admire Semitic fairy tales? You must be kidding! Anyone who reads the Koran cover to cover --is well--how shall I put it--not impressed--although one would see it as something the best a desert culture (with sand and sheep and no water) could produce, and in the typical high minded way just consider it ``different'' and not look down upon it unless it is imposed upon one--which is what the ``semitic'' traditions with their ``one book'' ``one way'' try to do upon others. Evola notes in the Doctrine of Awakening that this is not what the Aryan mentality does not try to ``impose'' in this fashion. Most of the great works in Islam are the work of people who were converted by the sword and forced into the religion, such as the Persians and Syrians and Lebanese, the latter two of whom were a part of the Greek colonies. As Dr. Sachau the great German orientalist (translator of Alberuni's Indica) writes: ``The foundations of Arabic literature was laid between AD 750 and 850. It is only the tradition relating to their religion and prophet and poetry that is peculiar to the Arabs; everything else is of foreign descent. The development of a large literature, with numerous ramifications, is chiefly the work of foreigners, carried out with foreign materials, as in Rome the origines of the national literature mostly point to the Greek sources. Greece, Persia, and India were taxed to help the sterility of the Arab mind... .. What India has contributed reached Baghdad by two different roads. Part has come directly in translations from the Sanskrit, part has travelled through Eran, having originally been translated from Sanskrit (Pali ? Prakrit ?) into Persian, and farther from Persian into Arabic. In this way, e.g. the fables of Kalila and Dimna have been communicated to the Arabs, and book on medicine, probably the famous Caraka.''

On the Arab knowledge of astronomy, Sachau writes:

``As Sindh was under the actual rule of Khalif Mansur (AD 753 -- 774), there came embassies from that part of India to Baghdad, and among them scholars, who brought along with them two books, the Brahamsiddhanta to Brahamgupta (Sirhind), and his Khandkhdyaka (Arkanda). With the help of these pandits, Alfazari, perhaps also Yakub ibn Tarik, translated them. Both works have been largely used, and have exercised a great influence. It was on this occasion that the Arabs first became acquainted with a scientific system of astronomy. They learned from Brahamgupta earlier than from Ptolemy.''

Also I find it hard to be impressed with Islam like Guenon. This man had a good understand of traditionalism and it is understandable that Guenon was impressed with Sufism. But what was Sufism really? The place where it arose-- Balkh-- was originally the confluence of Brahmanic, Buddhist and Zoroastrian civilizations before the inhabitants were sadly converted by the sword to Islam. After their conquest the ethnic composition of the peoples also changed. The consequences? All learning as vanished in the place Balkh ( today Afghanistan) which once produced a Rumi. Look at it today--their history begins with the birth of Mohammed. Around the time of Rumi, the original Brahmanic, Zoroastrian and Buddhist learning had not vanished and this is essentially where the ``mysticism'' in Islam comes from. To me, Sufism is the original Persian Aryan instinct rebelling against a matter of fact semitic religion imposed upon Persians via the sword. The religion was imposed upon them via force, so they added Persian mysticism to it. Read the conquest of Persia by the Arabs. It will make you weep. It took nearly a century for the Persians to be fully converted via the sword. Ask any educated Iranian and they will tell you that even though they were forcibly converted, they gave Islam a culture.

In Kashmir, Sufis practised ``takiya'' or deception--they imitated the Hindu and Buddhist practices so the inhabitants were drawn to Sufism, not really understanding the true nature of Islam and thinking that the Sufis were just another sect as in Hinduism and Buddhism--which they were not!!!

Read the famous explorer Richard Burton. He also states that most of the mysticism in Sufism is from the pre-Semitic influence i.e. Zoroastrian, Brahminism and Buddhism and Burton notes in his writings that this fact is not given due credit. Why? It is largely the nature of Semitic religions to think they have a monopoly on the truth--it is unAryan not to give credit where it is due... 

Guenon for all his understanding of initiation and tradition, does not understand Islam--its roots and origins. Anyone steeped in the older, original wisdom of mankind is not likely to be impressed by it. As Arthur Schopenhauer, perhaps the wisest of Germans, put it: '' the ancient wisdom of the human race will not be supplanted by the events in Galilee.''

And as for Christianity, it is easy to see that whatever is good in it was largely due to the fact that it became the religion of civilized people, i.e. Greeks and Romans and what is good in it can be traced to these sources and to later contributions when other barbarian lands in Europe became civilized and Christian.

So I fail to understand Evola's praise of Islam in some of his writings. Is it because the Nordics had nothing but Valhalla before being civilized by Greek and Roman civilization that he thinks that a violent religion which tires to uproot something superior and impose itself is something good as in the case of Islamic conquests of Central Asia? Evola has great insights into traditionalism, but sometimes he goes too far in trying to establish the ``Nordic'' myth. I do not think the old Indians and Iranians would think themselves related to the Europeans just as the ancient Greeks would not think themselves related to the barbarians surrounding them--the Gauls, Goths and what not... 

The ``Aryan'' concept has been rendered a lot of damage by white supremacists--in the East it is the civilizational ethos that counts--a Brahmin can be blue eyed (getting rarer with more descent to the plains), dark, or fair yet it is the civilizational ethos that is paramount (the dark and the fair have the same civilizational ethos--this is what matters ultimately--a point Evola also makes.) I hate the fact that the word ``Aryan'' which is used by Persians and Hindus, our ancestors, to refer to themselves, has been misappropriated by white supremacists and made to refer to refer to other peoples to whom the Hindus and Persians would not consider themselves to be related to--just as the Greeks referred to theselves as Hellene and non Greek speakers as barbarians, it was the same for Persians and Hindus--they referred to others as ``outsiders''. So when people ask are Indians `Aryans'' I laugh--its the Germanic peoples who were late comers into civilization (although they created the greatest modern one) who were never referred to as Aryans--guess the white supremacists wanted to extend their history to eras when they were still living in mud huts when the Greeks had civilization and relate themselves to the true ``Aryans'' to make themselves civilized all along and thereby feel superior and forget that once they were barbarian... .One has to see the mischief in Evola with his `Nordic'' myth and people like Guenon while appreciating their strong points which are many... .

It is sad for Europeans that they are Christians and cannot understand their pagan past apart from a serious study of Greek and Latin--but there are those of us in the East who are directly connceted to our distant ancestors having preserved the traditions and still practise the religion of our ancestors. A European needs to shed a lot to get out of his Christian mould and fully appreciate his true Greco-Roman heritage... For us I think we have an advantage in this respect--we are still very much connected to our Aryan past and do not need a ``Aryan'' mythology like white supremacists to feel superior... \hfill

\texttt{kadambari on 2010-06-19 at 22:52 said:}

Sorry for the typos above--forgot to check before submitting and there is no edit function... \hfill

\texttt{Cologero on 2010-06-20 at 00:18 said:}

Please bear in mind that the \emph{City of God} was written by a Christian monk. Therefore, for his ideal city to be populated by Indians instead of Europeans, and their religion based on Hinduism instead of Christianity, is quite significant. That is why I emphasized the point.

Similarly, for Mohammed, whom Dante has placed in Hell just a few centuries prior. The point is that by including the major figures in history, Campanella is anticipating Comte's later religion of Humanity.

Guenon's view of Islam and reason for his ``conversion'' is dealt with here: \textit{Did Guenon Convert?}\footnote{\url{http://www.gornahoor.net/?p=262}}

Guenon used the language of the Vedanta to express his metaphysical views. He gave as a reason for not ``converting'' to Hinduism is that it is too tied to the caste system, effectively rendering conversion for someone born outside it, basically impossible. Sufism was his second choice. Actually third, because for a time he thought he could influence Catholicism to take a more Traditional turn.

Evola often referred to Persian and Hindu mythology as ``Aryan'', and claimed ancient India as one of the ``three great Aryan civilizations'' (the others being the Greco-Roman and the Medieval). Evola used the terms Nordic and Aryan much more broadly than his epigones do todsy.

He regarded the Greeks and Romans as Nordics, and included some Asian and Central American peoples among the Aryans.

\hfill

\texttt{kadambari on 2010-06-20 at 06:52 said:}

``Evola often referred to Persian and Hindu mythology as ``Aryan'', and claimed ancient India as one of the ``three great Aryan civilizations'' (the others being the Greco-Roman and the Medieval). Evola used the terms Nordic and Aryan much more broadly than his epigones do todsy.''

He regarded the Greeks and Romans as Nordics, and included some Asian and Central American peoples among the Aryans.''

As for ``Nordic'' I can tell you my father in law was a tall blue eyed six foot Kashmiri Brahmin, extremely fair. But he would not think of himself as the same as European however much he looked like one. So it is the civilizational ethos that is carried which is of importance.

Yes I was merely pointing out the mischief the Aryan myth lead to--50 million dead in Europe from which Europe never recovered till today culturally speaking. Also I see quite a bit of mischief in Evola--his constant use of Sanskrit terms and concepts labelling them ``nordic'', when they have no relation to Nordic. This is just an example of how he tried to create a mythology for the Nordics whose real history began after being civilized by the Greco-Romans--does not matter who Evola considered ``Aryans''. This is my point. The Greeks, ancient Hindus and Persians would consider all people Evola labels Aryans as barbarian or outsiders. Today Greece is not the same--they became Christian and different--and whatever greatness the West has, its is obvious it has always been great when it has returned to its real roots--the Greco-Roman. I notice too much gloricication of Islam in Evola--and considering Islam often destroyed superior civilizations through brutality and force and left the people with a historical amnesia of their past such as in Central Asia, glorification of Islam is a mistake and the weakest point in his writings.

It is easy to see how Evola can be misleading for people who do not know their history too well--i.e. often take his mythology literally. One sees in the case of Nazis who were merely white supremacists what the Aryan myth lead to. To me Aryan means non-Semitic--in terms of philosophical and civilizational outlook--not the term it has come to connote due to the mischief of white supremacists.

So while I think Evola is an extremely intelligent writer, I also read between the lines and discard the nonsensical parts of his writings. He is amazing however in his understanding of tradition, the heirarchical nature of the Greco-Roman, Indic, Persian civilization etc. His weakest is his Nordic mythology. Physical resemblance does not entail cultural or philosophical resemblance... which is why when Greece was civilized the surrounding world around it in Europe was not... until they absorbed the civilization... .

\hfill

\texttt{kadambari on 2010-06-20 at 11:58 said:}

``He gave as a reason for not ``converting'' to Hinduism is that it is too tied to the caste system, effectively rendering conversion for someone born outside it, basically impossible.''

As regards the much maligned caste system, this is basically what kept India intact in the face of constant barbarian invasions by Afghan turks. Hindus became ever more conservative just to be able to survive Islam, and hence, caste gained in extreme rigidity in those areas where Islam was dominant. Native movements like Buddhism and Jainism had already begun to challenge caste and would have created a native internal solution to the problems presented by caste, but with the downfall of the Buddhist strongholds in Central Asia to the Muslims, this self-healing was stopped with the wholesale destruction of the Buddhist centers of learning. However, it is true that Hindus were often shortsighted as well. This has to do with being invaded constantly. In Kashmir, the Hindus converted to Islam wanted to convert back at one point, but the Brahmins were short sighted to allow for that. What I find amazing in India, is that there is something distinctly intact--although civilization has become sickly--despite all the semitization that occurred with the Turk invasions. Most of the North of India is a cesspool of crime and barbarism due to the destruction of the classical civilization in India by Muslims. In sections untouched by Muslims historically you find the people benign and fairly civilized.

Keep in mind however that were it not for caste despite its shortcomings there would be nothing left of anything good in India and India would definitely be Muslim like the central Asian brothers with complete historical amnesia of their past... .

\hfill

\texttt{Cologero on 2010-06-20 at 12:17 said:}

As you undoubtedly know, Evola ascribed race to a spiritual attitude, not to blood; of course, he was opposed to Nazi ideology in this regard. Furthermore, I think you are missing a certain perspective.

A healthy society knows what it is, hence what is not it. It is only a spiritually unhealthy culture that needs to go in search of its ``primordial Nordic tradition'', or ``Aryan'' heritage. So Evola's project is one of a rediscovery and reconstruction. Of course, he takes myths ``literally'', since they reveal the structure an ancient thought and social arrangements ... what he calls metahistory.

I don't know what ``\,''Aryan myth'' you claim the Nazi's adopted. Quite to the contrary, they adopted the Semitic myth of a ``chosen people' and applied it to themselves. Of course, there can be just one ``chosen people'', hence ... You may want to be more cautious about ``reading between the lines''; things may not always fit your preconceptions. Sometimes a cigar is just a cigar.

\hfill

\texttt{kadambari on 2010-06-20 at 13:01 said:}

``\,''Aryan myth'' you claim the Nazi's adopted''

well the Nazis misappropriated a lot of concepts of their scholars and used it for their agenda--they were not sophisticated people--just white surpemacists, but they did do a lot of damage with knowledge that they did not understand. One can hardly use the word ``Aryan'' anymore without all kinds of bad connotations and I get tired of reading of Sankrit passages attributed to the `Nordic way'' by Evola. Clearly he too has an agenda at times, although a very subtle one--and I don't begrudge him that, he was invloved in the politics of his time(but unsuccessful at it)--he tries to create a mythology or a history for a peoples at a time when they did not have a civilized history... thats what his Nordic solar path or whatever he calls it amounts to... I am sorry, those of us who are still heirs of a living ancient tradition do not like the works of our distant direct ancestors misappropriated for use by white supremacists... 

\hfill

\texttt{Matt on 2010-06-20 at 15:26 said:}

Kadambari

When Evola used the term Nordic, he most likely meant that which has its origins in the north. This is why he spoke of civilizations like Ancinet Persia, Vedic India, Inca empire, Aztec Empire, Rome, Greece, the Germanics etc. as bein Nordic in nature as their myths and ideas have a common view in that they saw their origins lying in the north of the world. In addition, besides the Incans and Aztecs, the Indo-Aryans, ancient Iranians, Romans, many of the Greek tribes, and the Germanic tribes and Celts are all Indo-European (what was called Aryan) having their common ethnic origin in a Proto-Indo-European people.

\hfill

\texttt{kadambari on 2010-06-20 at 17:57 said:}

Matt

I understand all of the above that you have mentioned--my study is in ancient languages: Greek, Latin etc. and some Sanskrit as well as European languages. However, you cannot deny that Nazis used the ideas of scholars to advance supremacist ideas in a narrow sense and did a lot of damage. If the Germans had been merely conservative they probably would have had the greatest civilization today--the high culture was destroyed by the excessive militarism... Also Aryan, swastika and such symbols have become tainted in a sense and interpreted in narrow fashion by lay people. This is what I am trying to say. I sense a civilizational commonality between the pagan civilizations: Persian, classical Indic, Greek, Roman etc. I think the core of these civilizations had ideals very different from the ideals arising in the Middle East which is the home of the three monotheistic religions--after all, there are two homes of the world religions: Middle East and India... However, Christianity took a great deal from Greece and Rome and moulded itself... I am a supporter of great civilizations arising anywhere because a great civilization confers benefits beyond itself... 

I am just saying that if the negatives were toned down a bit and less emphasis were laid on race in a purely biological sense, the Aryan ideals would appeal to a great number of people in a positive sense both in the East and West... I understand that the framework of the West is Christian... but it seems to be losing ground, and I doubt the Church will have as much influence as it did in the middle ages... It will be interesting to see how the West deals with coming up with alternatives to satisfy people. In the East, the loss of caste etc, does not really upset the essential teachings of the religions at the core... The problems are of a different kind: that posed by an artificially imposed secular state which is at odds with the native culture.

\hfill

\texttt{Cologero on 2010-06-20 at 22:34 said:}

Precisely. That is all outlined in Chapter 25 of \textbf{Revolt Against the Modern World}, \emph{The Northern-Atlantic Cycle}, so there is no need to rehash it all here. Evola accepts the Hyperborean origin of the Indo-Europeans. Actually, educated Hindus used to accept a similar story of the origins of Vedic civilization, though it is more or less not politically correct in our time.

\hfill

\texttt{Cologero on 2010-06-20 at 22:48 said:}

You are correct that the excesses of Nazism led to the dominance of leftism and anti-tradition. Although Evola emphasized the spiritual dimension of race, he left enough ambiguity for his followers today to focus excessively on race in a biological sense.

As far as the ``loss of caste'', both Evola and Guenon see that as a symptom of decline. The opposite of caste is egalitarianism, and that really does upset the essential teachings of any traditional religion.

\hfill

\texttt{kadambari on 2010-06-21 at 23:04 said:}

I don't know personally about loss of caste as we have not lost it in 3000 years, but I do know that being a Brahmin by birth is not necessarily the same as one by knowledge... and only in rare instances do the two coincide perfectly... Who says to lose caste is to believe in egalitarianism? There will always be inequality of minds... I am just saying that the fundamental principles of our religions does not change by loss of caste... 

\hfill

\texttt{kadambari on 2010-06-21 at 23:07 said:}

Why use Hindu symbols and Zoroastrian symbols? Why not dig into Nordic ones only? Why not Vandal Nordic solar tradition? Is it because there is not much ``tradition'' in it? Why bring in other people's culture to extend the history of the Nordic? A bit peculair to me.

\hfill

\texttt{Cologero on 2010-06-21 at 23:41 said:}

Perhaps I introduced a confusion. Caste is not necessarily passed on genetically, and probably should not be. So if you mean that caste by birth is disappearing in India, that may be true. But, as you pointed out, that does not preclude the existence of caste in some other form.

\hfill

\texttt{Cologero on 2010-06-21 at 23:49 said:}

As was pointed out in another comment, from Evola's perspective, the sources of Hindu and Zoroastrian symbols are themselves Nordic. He believed the Nordics originated in the polar regions and migrated in two waves: from North to South, and West to East. Therefore, the use of those mythologies is valid, in order to reconstruct the Nordic primordial tradition.

Yes, you are correct, the Germanic mythology is not identical to the primordial tradition. See, for example, \textit{Ultima Thule: Julius Evola and Herman Wirth}\footnote{\url{http://www.gornahoor.net/library/UltimaThule.pdf}}.

\hfill

\texttt{kadambari on 2010-06-23 at 17:36 said:}

Well I am not so interested in race that much as in the Aryan ideals--I think the Greeks, Persians, Hindus basically shared certain fundamental metaphysical ideals despite geographic barriers (these three being contemporary civilizations). The rest of Europe entered civilization much later--moreover, unlike the surviving Brahmins and Zoroastrians (almost extinct), no one in Europe has the continuity of a singular tradition that old--a historical, spiritual and racial memory arising from the cintinuity of a tradition... But then Greece bacame different after becoming Christian.. But as to the influence of race... well all I can say is that higher metaphysical and spiritual ideals is an Aryan phenomenon for sure--I find these ideals quite distinct from the ideals arising from the ideals of the three Mid-East monotheisms. I guess not being European, I am not interested in Nordic tradition--apart from what exists after Northern European became civilized through the spread of Greco-Roman civilization... .

\hfill

\texttt{Matt on 2010-06-25 at 14:16 said:}

Kadambari

I think its unfair to say that Northern European societies were not civilized until the spread of greco-roman civilization. By the time Rome was spreading itself into northern Europe, it had already degenerated into aphrodistic forms and by foreign ideas.

Sure, the northern europeans were not sophisticated as the roamsn in terms of technology and art, but they strongly held on to noble ideas like honour, fidelity community etc. that had significantly declined in Rome by that time. So one could say the northern europeans were somewhat barbaric in their lifestyle, though that could be said for all peoples who were nomadic (the founding tribes of vedic India) or nomadic descendants, but they certainly were not barbaric in their worldview or rituals.

\hfill

\end{sffamily}
\end{footnotesize}

